\section{Introduction}
An instruction to paint ``in the style of Monet'' can introduce a common
painting appearance, exaggerate familiar color or texture cues, or alter the
features that distinguish painters. These effects are easy to
conflate. Moving closer to Monet's collection, for example, does not show that
a generator reproduces the differences between Monet and Sisley: a generic
painting instruction could improve proximity to both.

We ask whether painter names induce differences aligned with those among
\emph{documented digital reproductions}, beyond a response shared across names.
For each painter, we compare its average color, spatial and texture measurements
with the four-painter average, then ask whether generated images reproduce
those labeled differences. The target belongs to this reference collection;
historical subject
mixtures and reproduction processes can contribute to its artist differences.
Holding the generated scene text fixed controls the prompt intervention; it
does not make those historical subjects identical.

The distinction matters for evaluation. A generator can produce easily
separable named conditions whose labels map to the wrong reference painters.
It can also match changes along the reference pattern while adding large
differences that do not follow it. Conversely, a restrained response may have lower squared
error than a stronger response, without distinguishing artists more accurately
in every respect. A single proximity score or recognition accuracy cannot
separate these possibilities.

Six model configurations receive the same scenes and prompt clauses.
Repeated outputs estimate
disagreement without counting random output variation as systematic error, under
the stated independence assumptions. Separate post-result checks examine
painter-pair differences, contrast magnitude, scene dependence and reference quality.

Our contribution is an empirical evaluation study with three findings: (i) the
shared measured response dominates total named-prompt change; (ii) positive
aggregate alignment can conceal weak distinctions between related painters;
and (iii) scene averaging, contrast rescaling and reference handling change what
model comparisons establish. Centering, cross-repeat distances and scalar
calibration are established operations; we use their joint interpretation to
show why response strength is insufficient evidence of labeled agreement.


\section{Related Work}
\citet{somepalli2024style} develop learned style descriptors and study similarity
across diffusion models. \citet{su2025prompted} evaluate prompted-artist
recognition using content-controlled name substitutions and compare named and
artist-free images with real-artist prototypes in CLIP space. Our question
concerns the labeled differences among reference painters after removing a
shared response; neither name substitution nor cross-model comparison is new.

\citet{deliege2025} assess stylistic imitation and stereotyping with experts.
AI-Pastiche studies authenticity judgments and prompt adherence
\citep{asperti2025pastiche}, while AI-WikiArt uses painting-derived captions and
vision-language models for attribution \citep{fu2025aiwikiart}. These tasks
distinguish convincing imitation, recognizable authorship and distributional
agreement. Here, the primary target is agreement among labeled painter-mean contrasts
in measured color, spatial and texture coordinates.

Quantitative art research studies color, composition and image statistics
\citep{kim2014,sigaki2018,lee2020}. \citet{kim2026} combine formal and contextual
representations for art-historical analysis and contextual image-to-image
experiments. Our text-only name intervention supplies no original composition.

Generative evaluation separates fidelity and diversity \citep{naeem2020},
while conditional metrics distinguish within-condition variation from relationships
among means \citep{benny2021conditional}. We use independent-partition inner products
as the primary noise-corrected squared-difference estimator \citep{diedrichsen2017representational}
and energy statistics as a complementary distributional measure \citep{szekely2013}.
Processing sensitivity \citep{parmar2022} and limits of learned-embedding
separation \citep{asperti2026} motivate the source-quality checks and the
restriction to a specified digital reference target.

\section{Design and Measurements}\label{sec:design}
\subsection{Painters and Reference Images}
The primary panel covers four related painters through 649 distinct reproductions: 297 Monet,
106 Sisley, 141 Pissarro and 105 C\'ezanne. Artist means receive equal weight,
so the larger Monet collection does not dominate the between-artist target.
A separate 221-work development panel sets common feature units; no primary
reference or generated image is used to fit that transformation. Both panels
were assembled before the present experiment and had been used in earlier
analyses. They are explicit, finite comparison panels rather than newly withheld
samples of each painter's complete oeuvre.

Reference files carry Wikimedia Commons source metadata and Wikidata work
identities \citep{commonsglam,vrandecic2014wikidata}. These support traceable
selection and rights records, but do not standardize photography, color
calibration or conservation state. These are measurements of digital
reproductions, whose capture conditions can affect the comparison.

\subsection{A Common-scene, Six-model Experiment}
We compare six requested service configurations, labeled GPT Image 1, GPT Image 2, GPT Image 2.5 Flare (Flare),
GPT Image 2.5 Sunburst (Sunburst), Nano Banana 2 and FLUX.2 Max (FLUX).
Fourteen fixed outdoor scene descriptions (briefs) span water,
built, land and mixed compositions. Each scene is crossed with six clauses:
no painting instruction; ``Render as an oil painting''; and ``Render as an oil
painting in the style of [painter]'' for each of the four painters. The scene
text is otherwise identical. Each model--scene--clause cell receives two
separately requested outputs, giving 1,008 images collected on 10 September 2026 (UTC).

All models receive text only and return 1,024-by-1,024-pixel images.
The four OpenAI requests specify medium quality; Nano Banana 2 specifies
1K resolution, while FLUX requests PNG output without an explicit resolution
field. These comparisons concern the recorded
configurations, not maximum attainable quality or equal expenditure per image.
Response records do not echo model identifiers, so these labels identify requested
configurations rather than independently verified checkpoints
(Appendix~\ref{app:provenance}).
The exact prompts, parameters and randomized request order were fixed before
collection. The 14 scenes are a fixed design,
not a random sample of all possible content.
Earlier, separate prompt interventions are summarized in Appendix~\ref{app:cohorts}.

Figure~\ref{fig:original-generated} shows all six conditions for one brief.
GPT Image 2's C\'ezanne trees have darker blue-green outlines than its Monet
and Sisley trees: a visible name response in this scene, whose agreement with
historical painter differences remains to be measured.

\begin{figure*}[p]
\centering\includegraphics[width=\linewidth]{specificity_original_generated.pdf}
\caption{One scene, six prompt conditions and six generators. The requested scene
is a broad river bending past a gravel bank, three trees on the opposite bank,
and a low hill under an overcast sky. In each model row, compare the four named
outputs with the artist-free and generic oil-painting controls on the right,
then compare the painter names with one another. The top row shows audited
reference regions for display; primary scores use full frames. These works were
neither generator inputs nor matched compositions. All generated examples use the first brief and first
repeat; Appendix~\ref{app:images} gives the full prompt, selection rule and provenance.}
\label{fig:original-generated}
\end{figure*}

\subsection{What the Features Measure}\label{sec:features}
The 31 features in Table~\ref{tab:features} capture properties that need not
change together. Two images can have similar chroma (color strength) but different
hue diversity, or similar total contrast but different edge directions and
fine-scale texture. Such distinctions make results by feature family interpretable.
Correlated coordinates and sensitivity to image processing nevertheless prevent
treating Euclidean distance as a validated measure of artistic style.



Images undergo orientation correction and conversion to sRGB using valid
embedded profiles; compatible unprofiled files are treated as sRGB. The primary
pipeline preserves aspect ratio at a 512-pixel short side, using Lanczos resizing
without upsampling. No painting-region mask is applied: source margins or
calibration strips can enter the measurements. Each coordinate is centered and scaled by its
median and interquartile range in the development panel, with equal total
weight per painter. Thus a unit is a development-panel IQR in that coordinate,
not a perceptual just-noticeable difference. The Euclidean metric weights
standardized coordinates equally, without further equalizing feature families.



\section{Comparing Painter Differences}\label{sec:method}
\subsection{What Counts as Agreement?}\label{sec:contrast-intuition}
The comparison has three steps. Measure the same 31 properties in every image.
For each historical painter, subtract the average of all four painters; this
records how that painter differs from the group. Repeat the subtraction among
the four named outputs for each generated scene, then compare the corresponding
painter differences. We call these differences \emph{contrasts}.

Consider median lightness: Monet averages $.325$ development-IQR units below
Sisley in the references. Across all 14 scenes and both repeats, GPT Image 2
has a smaller gap, $.254$, in the same direction. Subtracting each set's
four-painter mean preserves these gaps; a shared brightness shift cannot create
this distinction. The full comparison uses all 31 coordinates and four painters.

Two scores answer different questions. The \emph{aligned response}, $\beta$,
measures how far generated contrasts extend along the reference contrasts;
$\beta=1$ matches that component's strength. The \emph{error}, $D$, counts total
disagreement, including differences in other directions. Consider three idealized
generators: one makes no painter distinctions, one exactly matches the reference
contrasts, and one doubles them. Their expected errors are respectively 1, 0
and 1. A stronger response can therefore have more error, and a near-unit
aligned response can coexist with substantial mismatch elsewhere.

\subsection{Computing the Two Scores}
Let $\mu_a\in\R^{31}$ be painter $a$'s standardized reference mean and
$r_a=\mu_a-\frac14\sum_b\mu_b$ its contrast. The total squared size of these
contrasts, $H=\sum_a\|r_a\|^2$, normalizes both scores. Here $\|u\|^2$ sums squared
coordinates, $\langle u,v\rangle$ sums matching-coordinate products, and
hats mark estimates. For generated measurements $z_{msak}$, indices denote
model $m$, scene $s$, painter $a$ and repeat $k\in\{1,2\}$:
\begin{equation}
 d_{msak}=z_{msak}-\frac14\sum_b z_{msbk}.
 \label{eq:center}
\end{equation}
Centering removes shared additive shifts, but shared rescaling or mixing
of coordinates can still change the contrasts
(Appendix~\ref{app:score-details}). With $S=14$ scenes, the aligned response is
\begin{equation}
\begin{gathered}
\widehat\beta_{msk}=\frac{\sum_a\langle d_{msak},r_a\rangle}{H},\\
\widehat\beta_m=\frac1S\sum_s\frac{\widehat\beta_{ms1}+\widehat\beta_{ms2}}2.
\end{gathered}
\label{eq:beta}
\end{equation}
A positive value indicates an overall response along the reference differences;
it does not require agreement for each artist pair. Values above one indicate
exaggeration along that direction. To measure error, multiply the discrepancies
from the two repeats rather than squaring either repeat's discrepancy:
\begin{equation}
\begin{gathered}
\widehat D_{ms}=\frac1H\sum_a
\langle d_{msa1}-r_a,\ d_{msa2}-r_a\rangle,\\
\widehat D_m=\frac1S\sum_s\widehat D_{ms}.
\end{gathered}
\label{eq:distortion}
\end{equation}
Under stable conditional means and independent, mean-zero repeat errors,
noise that does not recur across requests contributes zero in expectation.
The expected score is then the normalized squared error of the scene-conditional
mean contrasts. Individual estimates can be negative; they are retained, not
interpreted as better-than-perfect recovery. The $D=1$ no-contrast benchmark
is a theoretical predictor, not the observed generic-painting arm.

The target is the same pooled historical contrast in every scene. A model
whose artist differences vary across scenes can therefore incur error even
if its across-scene average matches the references. Such variation need not
be an artistic mistake. The primary score audits whether the pooled reference
contrast recurs across the tested scenes; it does not identify the historically
appropriate contrast for each generated composition. We separately report mismatch after averaging scenes
and the added scene-variation component; their sum is $D$. Exact identities
and derivations are in Appendix~\ref{app:score-details}.

\subsection{Model Comparisons and Uncertainty}
We compare models scene by scene, then average their paired error differences.
This controls for the common brief; features, painters and individual images
are not treated as independent replicates.

Before inspecting new feature outcomes, we declared six aligned responses and
15 model-pair error differences. Paired-scene Student intervals use Bonferroni
adjustment for these 21 comparisons together, giving approximate 95\%
simultaneous coverage. Intervals excluding zero resolve their comparisons.
Absolute-$D$ intervals are nominal 95\% summaries: they are unadjusted and descriptive.

The intervals reflect variation across the 14 fixed briefs, combining generation
noise with real scene differences. They neither isolate fresh-request uncertainty for this exact
panel nor establish coverage for unseen scenes. Two outputs enable noise
correction but do not establish its sampling distribution or service independence.

Declared checks vary scenes, reference samples, feature families and measurement
windows. A separately declared sensitivity equalizes four title-derived subject
classes within each painter. These choices probe dependence on the reference
panel and metric; they do not match historical compositions. Their full
specifications remain in Appendix~\ref{app:estimation}.

\subsection{What the Additional Comparisons Test}\label{sec:diagnostics}
The artist-free and generic-painting arms characterize the common response removed
by centering.
Retrospective diagnostics address three questions: which painter pairs contribute
to the aggregate response, how contrast magnitude changes error, and how reference
choices change the comparison. For the magnitude check, we choose one
nonnegative multiplier per model, shared across all painter contrasts and
coordinates, to minimize error on 13 scenes.
We evaluate it on the omitted scene, rotating through all 14 scenes.
This \emph{calibration} multiplies only centered generated contrasts; reference
contrasts, coordinate units and images stay fixed. These diagnostics
interpret the declared results without extending their significance-test family.
Formulas appear in Appendix~\ref{app:score-details}; controls, fits and the source
audit are in Appendix~\ref{app:diagnostics}.

\section{Results}\label{sec:results}
\subsection{A Response in the Right Direction Can Still Miss the Target}
All 1,008 outputs yield valid measurements. With the full-frame
reference target, every model's aligned-response interval lies above zero
(Table~\ref{tab:specificity-models}). A purely common additive response is
therefore insufficient to explain the named conditions under the stated
inference model. Yet positive alignment does not ensure low error. GPT Image 2
nearly matches the reference-aligned strength, $\beta=.999\;[.893,1.104]$,
while its total error is $D=1.226$: it adds substantial differences beyond
that aligned component.

\begin{table*}[tbp]
\caption{Agreement with the full-frame reference target. $\beta=1$ matches its aligned amplitude; larger is not necessarily better. Expected $D$ is 0 for exact agreement and 1 for no painter distinctions. Positive $\Delta D$ favors FLUX.2 Max. Table brackets give approximate 95\% simultaneous intervals from the prespecified family of six aligned responses and 15 scene-paired model contrasts. FLUX's lowest $D$ estimate has an unadjusted 95\% interval $[0.586,1.016]$, spanning the no-distinction value 1; this interval is descriptive.}
\label{tab:specificity-models}

\centering\small\setlength{\tabcolsep}{4pt}
\begin{tabularx}{\linewidth}{@{}>{\raggedright\arraybackslash}p{.24\linewidth}>{\raggedright\arraybackslash}Xr>{\raggedright\arraybackslash}X@{}}
\toprule
Model & Aligned response $\beta$ & \shortstack{Uncalibrated\\error $D$} & $\Delta D$ versus FLUX\\
\midrule
GPT Image 1 & $0.938\;[0.709,1.168]$ & $1.572$ & $0.771\;[-0.037,1.580]$\\[3pt]
GPT Image 2 & $0.999\;[0.893,1.104]$ & $1.226$ & $0.425\;[-0.161,1.011]$\\[3pt]
GPT Image 2.5 Flare & $0.772\;[0.680,0.864]$ & $1.738$ & $0.937\;[0.256,1.618]$\\[3pt]
GPT Image 2.5 Sunburst & $0.824\;[0.680,0.968]$ & $1.641$ & $0.840\;[0.058,1.623]$\\[3pt]
Nano Banana 2 & $0.442\;[0.256,0.629]$ & $1.109$ & $0.309\;[-0.847,1.464]$\\[3pt]
FLUX.2 Max & $0.470\;[0.239,0.701]$ & $0.801$ & $0$ (baseline)\\[3pt]
\bottomrule\end{tabularx}

\end{table*}


FLUX.2 Max has the lowest error estimate, $.801$, despite a smaller aligned
response, $.470$. Adjusted comparisons resolve higher error for Flare and Sunburst
than for FLUX; the other 13 intervals include zero (Appendix~\ref{app:results},
Table~\ref{tab:all-specificity-pairs}).
These use full-frame references;
the Sunburst result changes after source correction
(Section~\ref{sec:source-correction}). No model is established to beat the
no-contrast benchmark of one. Thus the study detects a
reference-aligned response more clearly than it establishes low total mismatch.

Appendix~\ref{app:geometry} visualizes the same labeled contrasts in two
reference-fitted axes; pairwise results below use all 31 features.



\subsection{Most Scene-averaged Change Is Common to the Painter Names}\label{sec:shared-change}
After averaging scenes, named-minus-artist-free shifts comprise a four-name mean
and painter-specific departures. The mean includes the oil-painting clause and
any common naming response: 82.5--95.7\% of repeat-corrected squared change
(Table~\ref{tab:shared-change}). Averaging emphasizes consistent shifts;
splitting within scenes retains scene-specific variation and gives 88.6--97.2\%
(Appendix~\ref{app:estimation}). $\beta$ measures departures along the reference contrasts.

\begin{table*}[htbp]
\caption{Shares of scene-averaged, repeat-corrected squared change from artist-free to named prompts, including the painting clause. Shared is the four-name mean shift; between-name is departure from it. The lower rows give retrospective comparisons of the shared named shift with the generic oil-painting shift, using the same artist-free baseline. The noise-corrected cosine measures their directional agreement (1 means parallel); the ratio divides their repeat-corrected squared magnitudes, shared by generic.}
\label{tab:shared-change}
\centering\small\setlength{\tabcolsep}{3pt}
\begin{tabularx}{\linewidth}{@{}l*{6}{>{\centering\arraybackslash}X}@{}}
\toprule
& \shortstack{GPT Image\\1} & \shortstack{GPT Image\\2} & Flare & Sunburst & \shortstack{Nano\\Banana 2} & \shortstack{FLUX.2\\Max}\\\midrule
Shared (\%) & 82.5 & 85.8 & 84.8 & 83.9 & 88.1 & 95.7\\
Between-name (\%) & 17.5 & 14.2 & 15.2 & 16.1 & 11.9 & 4.3\\
\midrule
Corrected cosine & 0.689 & 0.774 & 0.766 & 0.837 & 0.930 & 0.916\\
Squared-size ratio & 1.986 & 1.960 & 2.816 & 3.356 & 3.692 & 4.138\\
\bottomrule\end{tabularx}

\end{table*}


The shared named shift points broadly along the generic oil-painting shift
(noise-corrected cosine $.689$--$.930$; 1 means parallel), but its repeat-corrected
squared magnitude is 1.96--4.14 times as large. Both use scene averages and the
same artist-free baseline; Table~\ref{tab:shared-change} lists each model's values.
Thus, a response shared across painter names need not be identical to the response
to a generic painting instruction. The generic shift and the additional shared
naming shift can reinforce or oppose one another, so their squared magnitudes
do not add as independent contributions (Appendix~\ref{app:estimation}).


\subsection{The Aggregate Response Conceals Uneven Painter Distinctions}
Retrospective pair diagnostics extend the lightness example in
Section~\ref{sec:contrast-intuition} to all 31 features (Figure~\ref{fig:artist-pairs}).
Monet--Sisley aligned responses range from $-.249$ to $.129$ in five models,
versus $.402$ for GPT Image 2. Yet GPT Image 1 has lower estimated error for this pair
($.56$ versus $1.16$): stronger alignment does not ensure lower mismatch.
Omitting C\'ezanne lowers FLUX's overall response from $.470$ to $.096$ and
GPT Image 2's from $.999$ to $.651$.

Swapping Monet and Sisley lowers error for Flare, Sunburst and Nano Banana 2.
Reversing a negatively aligned pair difference lowers error without changing its size.

\begin{figure*}[!htbp]
\centering\includegraphics[width=\linewidth]{specificity_artist_pairs.pdf}
\caption{Descriptive painter-pair estimates (no pairwise tests). M, S, P and C denote Monet, Sisley,
Pissarro and C\'ezanne. Pair aligned response 0 means no aligned component,
1 matches reference strength, values above 1 exaggerate it, and negative values
indicate the opposite direction.
Expected error is 0 for exact agreement and 1 for no distinction. Each pair uses
all 31 features and is normalized by its squared reference distance.
Negative error can arise from repeat correction,
not better-than-exact agreement.}
\label{fig:artist-pairs}
\end{figure*}


\subsection{Scene Averaging and Contrast Magnitude Change the Ordering}
We compare fit scene by scene, after averaging scenes, and after held-out
rescaling of painter contrasts (Table~\ref{tab:review-calibration}).

\emph{Averaging scenes.} Nano Banana 2's mean painter contrasts have lower error than FLUX's
($.723$ versus $.761$). Evaluating each scene adds a larger variation component
for Nano Banana 2 ($.387$ versus $.040$), reversing that order.
$D$ penalizes variation around the pooled historical target; averaging can
cancel scene-dependent differences.

\begin{table*}[!htbp]
\caption{Retrospective error decomposition and calibration against full-frame references. $D=D_{\rm aggregate}+V_{\rm scene}$. Calibration rescales centered generated contrasts using the other 13 scenes and evaluates them on the omitted scene; it does not produce new images. Bold marks each error column's minimum, not significance.}
\label{tab:review-calibration}
\centering\small
\begin{tabularx}{\linewidth}{@{}l*{3}{>{\centering\arraybackslash}X}@{\hspace{2em}}>{\centering\arraybackslash}X@{}}\toprule
Model & \shortstack{Uncalibrated\\scene error\\$D$} & \shortstack{Error after\\averaging scenes\\$D_{\rm aggregate}$} & \shortstack{Scene\\variation\\$V_{\rm scene}$} & \shortstack{Calibrated\\scene error\\$D_{\rm held}$}\\
\cmidrule(r){1-4}\cmidrule(l){5-5}
GPT Image 1 & 1.572 & 1.239 & 0.333 & 0.645\\
GPT Image 2 & 1.226 & 1.044 & 0.182 & \textbf{0.554}\\
2.5 Flare & 1.738 & 1.483 & 0.255 & 0.741\\
2.5 Sunburst & 1.641 & 1.337 & 0.305 & 0.706\\
Nano Banana 2 & 1.109 & \textbf{0.723} & 0.387 & 0.832\\
FLUX.2 Max & \textbf{0.801} & 0.761 & 0.040 & 0.714\\
\bottomrule\end{tabularx}

\end{table*}


\emph{Rescaling contrasts.} GPT Image 2's aligned response is near one, but its
squared contrast size exceeds twice the reference's ($Q=2.223$, where 1 matches the reference size;
Appendix~\ref{app:magnitude-coverage}).
Shrinking weakens its reference-aligned component, but reduces off-target
differences enough to lower the total error estimate. GPT Image 2's fitted multipliers
scale every contrast coordinate by about 0.45, using only the other 13 scenes each time.

After calibration, GPT Image 2 has the lowest error estimate, $.554$,
compared with FLUX's $.714$; it has lower scores in 12 of 14 held-out folds.
Deleting any scene and refitting the entire procedure preserves the mean
advantage (Appendix~\ref{app:magnitude-coverage}).
The folds overlap, so their win count does not establish statistical superiority.

\subsection{Which Comparisons Survive Reference Changes?}\label{sec:source-correction}
For the uncalibrated scene error $D$, FLUX retains the lowest point
estimate across scene deletions, measurement windows, reference targets and
feature-weighting sensitivities (Appendices~\ref{app:declared-sensitivities}
and~\ref{app:diagnostics}).

Source handling has a more consequential effect on statistical comparisons.
AI assistants inspected all reference and development reproductions, selecting
crops that exclude peripheral artifacts in 90 reference and 41 development images.
The audit used no human verification. After cropping and refitting development scaling,
FLUX's error estimate changes from $.801$ to $.872$. Corrected aligned-response
intervals remain above zero for all six models, and GPT Image 2 retains the lowest
calibrated estimate.

Table~\ref{tab:source-comparison} makes the changed comparisons explicit. With
corrected regions and refitted scaling, the Sunburst interval crosses zero,
Flare still has higher error than FLUX and the GPT Image 1 interval moves above zero. The latter
is a retrospective sensitivity, not a new confirmatory finding. Source correction
preserves FLUX's lowest point estimate while changing the evidence for particular
model differences.
Weak Monet--Sisley alignment also persists after correction, but individual
signs are unstable: Flare changes from $-.011$ to $.019$, GPT Image 1 from
$.074$ to $-.053$, and FLUX from $.129$ to $.019$. The label-swap observation
above therefore concerns the original reference target.


\begin{table*}[!htbp]
\caption{Uncalibrated error differences before and after source correction, using audited painting regions and refitted development scaling. Positive differences favor FLUX.2 Max. Full-frame intervals use the 95\% simultaneous procedure across 21 endpoints; corrected intervals reuse that procedure retrospectively and add no confirmatory tests.}
\label{tab:source-comparison}
\centering\small
\begin{tabularx}{\linewidth}{@{}Xcc@{}}\toprule
Model comparison & \shortstack{Full-frame difference\\{[adjusted interval]}} & \shortstack{Crop + scaling correction\\Difference [adjusted interval]}\\\midrule
GPT Image 1 $-$ FLUX & $0.771\;[-0.037,1.580]$ & $0.758\;[0.026,1.491]$\\[2pt]
2.5 Flare $-$ FLUX & $0.937\;[0.256,1.618]$ & $0.864\;[0.193,1.535]$\\[2pt]
2.5 Sunburst $-$ FLUX & $0.840\;[0.058,1.623]$ & $0.705\;[-0.026,1.436]$\\[2pt]
\bottomrule\end{tabularx}

\end{table*}



A separate check uses class-specific reference means for 11 water, built and
land briefs, excluding three mixed briefs. Replacing title-derived reference
labels with visual labels retains FLUX's lowest error estimate.
Appendix~\ref{app:diagnostics} gives the controls and audit.\par



\section{Discussion}\label{sec:discussion}
Detecting a name response and matching painter differences are distinct achievements.
FLUX's lowest uncalibrated error estimate coexists with weak Monet--Sisley alignment.
Model comparisons should therefore report overall and pairwise scores together
for the stated reference target, distinguishing estimates from resolved comparisons.
Reports should also state whether they score each scene, the average pattern,
or held-out contrasts rescaled using other scenes: these choices change the model ordering.

The intervention identifies differences between prompt conditions, not their
internal causes. Training composition, guidance and shared transformations remain
possible explanations.

\subsection{Scope and Remaining Limitations}
The target is the pattern of differences in a finite digital reference
collection. It mixes subject choices and reproduction processes with properties
of the paintings. Coarse content classes do not match composition, viewpoint,
career period or actual generated scene adherence. The source audit addresses
visible peripheral regions and title/visual-class disagreements, but AI
judgments, uncertain boundaries and remaining capture differences limit that
correction. One reproduction per work cannot separate effects of illumination,
sharpening, conservation or digitization source.

The 31 coordinates describe digital color, spatial organization and texture,
not physical brushwork or perceived painter resemblance; numerical replay and
robustness to coordinate weights do not validate those judgments.
Inference also conditions on four related
painters, 14 authored briefs, two repeats and one clause template. Simulations
probe stated noise models, not actual service independence. The image panels
make the comparison inspectable, while incomplete public exact-pixel access
still limits independent measurement reproduction.

\section{Conclusion}
Most squared feature change from artist-free to named prompts is shared across
names after scene averaging, including the oil-painting contribution.
All six models nevertheless show reference-aligned painter differences, yet none
is established to beat the no-distinction error benchmark before calibration.
Uneven painter-pair agreement and the changed error ordering after calibration
show why a single overall score is insufficient.
